\section{模型设计与评估方案}

\subsection{总体技术流程}

本项目以一条清晰主线组织完整分析流程：

\begin{center}
\begin{tabular}{c}
\hline
\textbf{跨年度字段和数据质量审计} \\
$\downarrow$ \\
\textbf{按年份划分训练、验证和测试集} \\
$\downarrow$ \\
\textbf{构造无数据泄漏的建筑属性特征} \\
$\downarrow$ \\
\textbf{建立建筑能耗基准回归模型（多模型对比）} \\
$\downarrow$ \\
\textbf{时间外测试（2023、2024）和未见建筑测试（GroupShuffleSplit）} \\
$\downarrow$ \\
\textbf{SHAP解释重要预测特征} \\
$\downarrow$ \\
\textbf{基于标准化残差和同类基准识别异常高耗能建筑} \\
$\downarrow$ \\
\textbf{按OSEBuildingID汇总最近三年表现为建筑级记录} \\
$\downarrow$ \\
\textbf{构建建筑级改造优先级评分并检验排序稳定性} \\
$\downarrow$ \\
\textbf{对高优先级候选建筑进行辅助聚类画像} \\
\hline
\end{tabular}
\end{center}

\subsection{特征工程}

\subsubsection{推荐输入特征}

从原始字段中构造以下特征用于能耗基准预测：

\begin{table}[H]
\centering
\caption{模型特征列表}
\label{tab:features}
\begin{tabular}{p{5cm}p{7cm}}
\hline
\textbf{特征} & \textbf{计算方式与说明} \\
\hline
BuildingType & 建筑类型，使用OneHotEncoder编码 \\
\hline
LargestPropertyUseType & 主要用途类型，使用OneHotEncoder编码 \\
\hline
PropertyGFATotal & 总建筑面积（需重新核查该字段在各年份的实际覆盖率，不应直接判定为全部缺失） \\
\hline
NumberofFloors & 楼层数 \\
\hline
NumberofBuildings & 建筑数量 \\
\hline
YearBuilt & 建造年份（异常值标记后按建筑类型中位数填补） \\
\hline
BuildingAge & $\max(\text{DataYear} - \text{YearBuilt}, 0)$ \\
\hline
Neighborhood & 所在区域 \\
\hline
ParkingRatio & $\text{PropertyGFAParking} / \text{GFA}$ \\
\hline
AreaPerFloor & $\text{GFA} / \text{NumberofFloors}$ \\
\hline
AreaPerBuilding & $\text{GFA} / \text{NumberofBuildings}$ \\
\hline
LogGFA & $\log(1 + \text{GFA})$ \\
\hline
IsMixedUse & 建筑用途类型数 $> 1$ 的指示变量 \\
\hline
PrimaryUseAreaRatio & 主要用途面积占GFA的比例 \\
\hline
DataYear & 数据记录年份（作为年度背景变量） \\
\hline
ES\_Missing & ENERGYSTARScore是否缺失的二值指示变量 \\
\hline
YearBuilt\_Anomaly & YearBuilt异常的指示变量（YearBuilt $>$ DataYear 标记） \\
\hline
\end{tabular}
\end{table}

\subsubsection{谨慎使用或剔除的字段}

当目标为SiteEUI时，主模型不得使用以下字段：
\begin{itemize}
    \item 当年SiteEnergyUse、SourceEUI、天气标准化EUI（由目标直接派生或换算）；
    \item 当年电力、天然气、蒸汽消耗量（与目标共享年度能耗测量信息）；
    \item 当年TotalGHGEmissions、GHGEmissionsIntensity（由能耗乘以排放因子计算，预测时不可预先获得）；
    \item 任何与目标存在确定性数量关系的字段。
\end{itemize}

GHG排放强度仅在优先级排序阶段作为减碳维度参考。

\subsubsection{ENERGY STAR Score处理}

设置两组实验：
\begin{itemize}
    \item \textbf{模型A（主结果）：}不使用ENERGY STAR Score，仅保留ES\_Missing指示变量；
    \item \textbf{模型B（扩展对比）：}使用ENERGY STAR Score作为特征，仅用于对比模型A，检验评分信息对预测性能的提升幅度。
\end{itemize}

\subsubsection{类别编码策略}

不同模型采用不同的编码方案，禁止对建筑类型使用LabelEncoder：
\begin{itemize}
    \item \textbf{ElasticNet / 线性回归：}OneHotEncoder + RobustScaler；
    \item \textbf{随机森林 / XGBoost：}缺失填补 + OneHotEncoder；
    \item \textbf{LightGBM：}可OneHotEncoder，也可使用原生类别特征（dtype="category"）；
    \item \textbf{K-Means聚类：}仅使用连续数值特征，建筑类型不参与距离计算，仅用于聚类结果的事后解释。
\end{itemize}

\subsection{训练、验证与测试划分}

\subsubsection{主实验：按完整年份扩展窗口}

\begin{table}[H]
\centering
\caption{按年份滚动验证方案}
\label{tab:split}
\begin{tabular}{p{3cm}p{3cm}p{2.5cm}}
\hline
\textbf{训练集} & \textbf{验证集} & \textbf{用途} \\
\hline
2015--2017 & 2018 & 交叉验证折1 \\
2015--2018 & 2019 & 交叉验证折2 \\
2015--2019 & 2020 & 交叉验证折3 \\
2015--2020 & 2021 & 交叉验证折4 \\
2015--2021 & 2022 & 交叉验证折5（模型选择） \\
\hline
2015--2021 & 2023 & 测试（时间外评估） \\
2015--2021 & 2024 & 测试（时间外评估） \\
\hline
\end{tabular}
\end{table}

最终在2023和2024两个测试集上分别报告性能，再汇总2023--2024结果。所有数据预处理（填充器、编码器、标准化器）仅在训练集上拟合。

\subsubsection{辅助实验：按建筑ID分组}

增加按OSEBuildingID进行的GroupShuffleSplit实验，训练集与测试集中的建筑完全不重叠，用于检验模型对``从未见过的新建筑''的泛化能力。

\subsection{模型选择与评估方法}

\subsubsection{模型序列}

模型按复杂度递增顺序评估：

\begin{enumerate}
    \item \textbf{基线A（全体均值）：}以训练集SiteEUI均值作为所有建筑的预测值；
    \item \textbf{基线B（建筑类型中位数）：}按建筑类型分组，以训练集各组SiteEUI中位数作为同类建筑的预测值；
    \item \textbf{ElasticNet：}带L1/L2正则化的线性模型，使用OneHotEncoder编码类别特征，随机搜索调参；
    \item \textbf{随机森林：}Bagging集成，缺失填补后OneHotEncoder编码，随机搜索调参；
    \item \textbf{XGBoost：}梯度提升树，OneHotEncoder编码，随机搜索调参；
    \item \textbf{LightGBM：}梯度提升树，支持原生类别特征处理，随机搜索调参。
\end{enumerate}

所有模型训练时对目标取$\log(1+\text{SiteEUI})$以缓解右偏分布的影响，预测后通过$\exp(\hat{y})-1$还原为原始尺度。

\subsubsection{评价指标}

在验证集与测试集上分别计算以下指标（均在原始尺度上报告）：
\begin{itemize}
    \item MAE（平均绝对误差）；
    \item RMSE（均方根误差）；
    \item $R^2$（决定系数）；
    \item 可选RMSLE（均方根对数误差）。
\end{itemize}

不将$R^2>0.75$设为刚性成功标准。更合理的目标是：最佳模型在测试集上的MAE应明显优于全体均值基线和建筑类型中位数基线（例如降低10\%--20\%）。

\subsubsection{分组报告}

评估时应分别报告以下子集的性能：
\begin{itemize}
    \item 2023年测试集指标；
    \item 2024年测试集指标；
    \item 不同建筑类型（办公楼、酒店、医院等）的指标；
    \item 高能耗建筑子集（SiteEUI高于同类型P75）的指标；
    \item 已见建筑（训练集中出现过的OSEBuildingID）与未见建筑的指标。
\end{itemize}

\subsection{SHAP可解释性分析}

仅对最佳模型和测试集进行SHAP分析：
\begin{enumerate}
    \item 计算测试集上的SHAP值；
    \item 输出SHAP全局重要性条形图；
    \item 绘制SHAP summary beeswarm图；
    \item 绘制BuildingAge与LogGFA的SHAP依赖图；
    \item 绘制典型高能耗建筑与低能耗建筑的对比瀑布图；
    \item 对高度相关特征进行分组分析，避免重复计数；
    \item 报告中统一使用``重要预测特征''表述，明确标注SHAP反映的是模型内部的关联模式，不等于因果关系，不使用``驱动因素''或``因果因素''等措辞。
\end{enumerate}

\subsection{异常高耗能建筑识别}

不再使用``残差$>0$''直接判定异常，改为三层递进筛选：

\subsubsection{第一层：预测残差异常}

计算各建筑的预测残差：
\[\text{Residual} = \text{ActualEUI} - \text{PredictedEUI}\]
再按建筑类型和年度进行标准化：
\[\text{StandardizedResidual} = \frac{\text{Residual} - \mu_{\text{type, year}}}{\sigma_{\text{type, year}}}\]
异常条件：标准化残差超过合理的统计阈值（如1.5倍标准差），或残差位于同类型同年度前10\%。

\subsubsection{第二层：同类基准超标}

计算相对于同类P75基准的超耗比例：
\[\text{BenchmarkExcessRatio} = \frac{\text{ActualEUI} - \text{BenchmarkP75}}{\text{BenchmarkP75}}\]
要求同时满足：超出同类型P75，且超耗比例达到合理分位数阈值。

\subsubsection{第三层：持续性验证}

检查2022--2024年中的持续表现：
\begin{itemize}
    \item 进入同类高耗能前25\%的年数；
    \item 残差进入同类前10\%的年数。
\end{itemize}
最终候选需同时满足三层条件：模型预测异常 + 同类基准超标 + 最近三年持续出现。

\subsection{建筑级记录聚合}

原始数据为``建筑—年份''粒度，最终需要每栋建筑一条优先级记录。按OSEBuildingID汇总：
\begin{itemize}
    \item \textbf{当前状态（2024年）：}ActualEUI、PredictedEUI、StandardizedResidual、BenchmarkExcessRatio、GHGIntensity；
    \item \textbf{持续性表现（2022--2024年）：}最近三年高耗能次数、最近三年平均残差；
    \item \textbf{建筑属性：}建筑名称、建筑类型、BuildingAge、GFA；
    \item \textbf{数据质量：}最近三年数据完整率。
\end{itemize}

每栋建筑最终仅保留一行。

\subsection{优先级评分体系}

\subsubsection{评分维度设计}

\begin{table}[H]
\centering
\caption{节能改造优先级评分维度}
\label{tab:scoring}
\begin{tabular}{p{2.5cm}p{1.2cm}p{2cm}p{6cm}}
\hline
\textbf{维度} & \textbf{权重} & \textbf{方向} & \textbf{计算方式} \\
\hline
标准化预测残差 & 25\% & 正向 & 同建筑类型、同年度的标准化残差，归一化到0--1 \\
\hline
同类基准超耗程度 & 20\% & 正向 & 超出同类型P75基准的比例（BenchmarkExcessRatio），归一化 \\
\hline
持续高耗能频率 & 15\% & 正向 & 2022--2024年中进入同类高耗能前25\%的年数比例 \\
\hline
理论超额能耗量 & 15\% & 正向 & $\max(\text{超额EUI}, 0) \times \text{建筑面积}$，归一化（不称为``预计节能潜力''，因未考虑改造成本、技术方案、实际节能率和投资回收期） \\
\hline
GHG排放强度 & 15\% & 正向 & 单位面积碳排放（GHGEmissionsIntensity），归一化 \\
\hline
建筑年龄 & 10\% & 正向 & BuildingAge标准化到0--1 \\
\hline
\end{tabular}
\end{table}

\subsubsection{数据质量折扣}

定义数据质量得分为：
\[\text{DataQualityScore} = 1 - \text{MissingRate}\]
最终得分：
\[\text{FinalScore} = \text{RawPriorityScore} \times \text{DataQualityScore}\]
数据缺失严重的建筑不会因某个极端值被错误排到前面。

\subsubsection{敏感性分析}

设置四组权重方案进行对比：

\begin{table}[H]
\centering
\caption{权重敏感性分析方案}
\label{tab:sensitivity}
\begin{tabular}{p{2.5cm}p{1.8cm}p{1.8cm}p{1.8cm}p{1.8cm}p{1.8cm}p{1.8cm}}
\hline
\textbf{方案} & 标准化残差 & 基准超耗 & 持续频率 & 超额能耗 & GHG强度 & 建筑年龄 \\
\hline
标准方案 & 25\% & 20\% & 15\% & 15\% & 15\% & 10\% \\
节能优先 & 35\% & 25\% & 15\% & 15\% & 5\% & 5\% \\
减碳优先 & 15\% & 10\% & 10\% & 15\% & 40\% & 10\% \\
大型建筑优先 & 20\% & 15\% & 10\% & 30\% & 15\% & 10\% \\
\hline
\end{tabular}
\end{table}

\subsubsection{排序稳定性验证}

由于没有真实改造结果标签，通过以下方式间接验证排序的稳健性：
\begin{itemize}
    \item 不同权重方案下Top-50候选清单的重合率与Jaccard相似度；
    \item Spearman秩相关系数与Kendall秩相关系数；
    \item Bootstrap重采样后每栋建筑进入Top-50的频率；
    \item 各建筑类型在Top-50中的占比是否合理（避免某类型垄断）；
    \item 与官方benchmark高耗能建筑列表的一致性检验。
\end{itemize}

\subsection{辅助聚类画像}

聚类不再负责判断``谁异常''（该任务已由三层异常识别完成），而是对高优先级候选建筑进行画像，辅助理解候选建筑的类型结构：

\begin{itemize}
    \item \textbf{聚类对象：}优先级排名较高的候选建筑（非全部建筑）；
    \item \textbf{聚类特征（仅连续变量）：}LogGFA、BuildingAge、BenchmarkExcessRatio、持续高耗能频率、GHGIntensity、理论超额能耗量；
    \item \textbf{建筑类型角色：}不参与K-Means距离计算，仅在聚类后用于解释各聚类的类型构成；
    \item \textbf{K值选择：}综合肘部法则、轮廓系数、Calinski-Harabasz指数确定；
    \item \textbf{画像命名示例：}``大型老旧高超耗办公楼''、``高排放高强度酒店''、``面积较小但连续异常的商业建筑''、``大体量高绝对节能规模建筑''。
\end{itemize}
